\section{CAMELYON17 natural center-shift studies}

CAMELYON17/WILDS \cite{bandi2018camelyon17,koh2021wilds} provides a complementary institutional setting in which center identity is observed rather than synthetically assigned. The feature-level dataset used here contains 455,954 labeled examples from five centers. Centers 0, 3, and 4 form the source-training pool, center 1 is used for out-of-distribution validation, and center 2 is held out for final testing.

\subsection{Source weighting and held-out-center generalization}
Three source-influence policies are compared: sample-proportional weighting, equal-client weighting, and downweighting the dominant source. The experiment is implemented as a centralized feature-level proxy for institutional influence, not as a full federated training run.

\begin{table}[t]
\centering
\caption{Held-out center-2 accuracy under two feature families and three source-weighting policies.}
\label{tab:camelyon_weighting}
\small
\begin{tabular}{lrrr}
\toprule
Feature extractor & Sample-proportional & Equal-client & Downweight dominant \\
\midrule
Frozen ImageNet ResNet18 & 0.8312 & \textbf{0.9132} & 0.9094 \\
CAMELYON17-trained ResNet18 & 0.9052 & 0.9318 & \textbf{0.9322} \\
\bottomrule
\end{tabular}
\end{table}

On frozen ImageNet features, equal-client weighting improves held-out-center accuracy by 8.20 percentage points over sample-proportional weighting. With source-trained CAMELYON17 features, the gap narrows but remains positive: dominant-source downweighting reaches 0.9322 versus 0.9052 for sample-proportional weighting. The source-trained sample-proportional model reaches approximately 0.9991 source-training accuracy, demonstrating that stronger source fit does not guarantee stronger performance on the unseen center.

These experiments reinforce the PANDA stress result from a different direction. More source samples can improve estimation when source and target distributions align, but sample volume alone does not determine which weighting best generalizes to a new center.

\subsection{Center subspace diagnostics}
The same dataset is used to study whether center identity occupies a partially removable representation subspace. Adversarial and subtractive nuisance-removal variants did not materially reduce post-hoc center leakage. A supervised linear center-subspace projection provided a cleaner mechanism test: five-center decodability fell from approximately 0.895 to 0.764 at effective rank four while tumor AUC remained approximately 0.990. This is evidence that part of the center signal is linearly structured and removable without a large loss on the measured tumor endpoint; it is not evidence of complete center invariance.

\subsection{Communication and systems studies}
The CAMELYON17 feature/head setting also provides a bounded systems calculation for federated communication. Across 100 fp32 rounds and three clients, exchanging a full ResNet18 was estimated at approximately 24.98 GB, compared with approximately 2.35 MB for a 512-to-2 classification head---a ratio of roughly 10,894. A coefficient-noise perturbation study preserved the equal-client and dominant-source-downweighting advantages through noise standard deviation 0.20, while a separate infrastructure simulation measured client-speed heterogeneity, dropout, failed rounds, and synchronous straggler burden.

These experiments are useful engineering bounds around the PA-NF institutional stage. They do not substitute for a real multi-hospital federated deployment, but they connect source weighting, center shift, communication cost, and infrastructure heterogeneity within the same research pipeline.
